\subsection{李群对流形的商}

设 \(\mathbf{G}\) 为李群，\(\mathbf{X}\) 为 \(\mathbf{C}^r\) 类流形，且 \((g,x)\mapsto gx\) 是 \(\mathbf{C}^r\) 类的 \(\mathbf{G}\) 在 \(\mathbf{X}\) 上的左作用律（Algebra, Chapter I, § 5, no. 1）。对于所有 \(g\in G\)，令 \(\tau (g)\) 表示自同构 \(x\mapsto gx\) 在 \(\mathbf{X}\) 上的映射，它由 \(g\) 定义。对于所有 \(x\in X\)，令 \(\rho (x)\) 表示轨道映射 \(g\mapsto gx\)，它从 \(\mathbf{G}\) 到 \(\mathbf{X}\)，由 \(x\) 定义。于是

\[
\rho (x) = \rho (g x) \circ \delta (g), \quad \rho (x) = \tau (g) \circ \rho (x) \circ \gamma (g ^ {- 1})\tag{3}
\]

对所有 \(g\in G\) 和 \(x\in X\) 成立。因此

\begin{enumerate}
\def\labelenumi{(\arabic{enumi})}
\setcounter{enumi}{3}
\tightlist
\item
\end{enumerate}

\[
\mathrm{T} _ {g} (\rho (x)) = \mathrm{T} _ {e} (\rho (g x)) \circ \mathrm{T} _ {g} (\delta (g))
\]

\[
\mathrm{T} _ {g} (\rho (x)) = \mathrm{T} _ {x} (\tau (g)) \circ \mathrm{T} _ {e} (\rho (x)) \circ \mathrm{T} _ {g} (\gamma (g ^ {- 1}))\tag{5}
\]

命题 9. 设 \(x \in X\)，\(g_0 \in G\)。

\begin{enumerate}
\def\labelenumi{(\roman{enumi})}
\item
  如果 \(\rho(x)\) 在 \(g_0\) 处是浸入（相应地是浸没、次浸入），则对于所有 \(g \in G\)，\(\rho(gx)\) 是浸入（相应地是浸没、次浸入）。
\item
  如果 \(\rho(x)\) 的秩为 \(k\)（在 \(g_0\) 处），则对于所有 \(g \in G\)，\(\rho(gx)\) 的秩恒等于 \(k\)。
\end{enumerate}

这立即由公式 (4) 和 (5) 得出，因为 \(\mathrm{T}_{g}(\delta(g)),\mathrm{T}_{x}(\tau(g))\) 和 \(\mathrm{T}_{g}(\gamma(g^{-1}))\) 都是同构。

推论。设 \(x \in \mathbf{X}\)。如果 \(\mathbf{K}\) 的特征为 0 且 \(\mathbf{X}\) 是有限维的，则 \(\rho(x)\) 是次浸入。如果进一步 \(\rho(x)\) 是单射，则 \(\rho(x)\) 是浸入。

这由命题 9 和 Differentiable and Analytic Manifolds, R, 5.10.6 得出。

注意，如果 \(\eta\) 表示映射 \((g,x)\mapsto gx\)，它从 \(\mathbf{G}\times \mathbf{X}\) 到 \(\mathbf{X}\)，那么对于 \(g\in \mathbf{G},x\in \mathbf{X},u\in \mathrm{T}_g(\mathrm{G}),v\in \mathrm{T}_x(\mathrm{X})\)，

\[
\mathrm{T} _ {(g, x)} (\eta) (u, v) = \mathrm{T} _ {(g, x)} (\eta) (u, 0) + \mathrm{T} _ {(g, x)} (\eta) (0, x)
\]

即

\[
\mathrm{T} _ {(g, x)} (\eta) (u, v) = \mathrm{T} _ {g} (\rho (x)) u + \mathrm{T} _ {x} (\tau (g)) v.\tag{6}
\]

命题 10. 设 G 为李群，X 为 \(\mathbf{C}^r\) 类流形，且 G 在 X 上有一个 \(\mathbf{C}^r\) 类左作用律。设：

\begin{enumerate}
\def\labelenumi{(\alph{enumi})}
\item
  群 G 在 X 上适当且自由地作用；
\item
  对所有 \(x \in \mathbf{X}\)，\(\rho(x)\) 都是浸入（如果 \(\mathbf{K}\) 的特征为 0 且 \(\mathbf{X}\) 是有限维的，这是 (a) 的推论）。
\end{enumerate}

G 在 X 上定义的等价关系是正规的（Differentiable and Analytic Manifolds, R, 5.9.5）。在商集 X/G 上存在唯一的流形结构，使典范映射 \(\pi: \mathrm{X} \to \mathrm{X} / \mathrm{G}\) 是浸没。该流形结构的底层拓扑是 X 上拓扑的商拓扑；它是豪斯多夫的。最后，(X, G, X/G, \(\pi\) ) 是主左纤维丛。 \(\dagger\)

令 \(\theta\) 为映射 \((g, x) \mapsto (x, gx)\)，它从 \(\mathbf{G} \times \mathbf{X}\) 到 \(\mathbf{X} \times \mathbf{X}\)。该映射是 \(\mathbf{C}^r\) 类的。我们证明它是浸入。对于 \(u \in \mathrm{T}_g(\mathbf{G})\) 和 \(v \in \mathrm{T}_x(\mathbf{X})\)，由 (6) 得

\[
\mathrm{T} _ {(g, x)} (\theta) (u, v) = (v, \mathrm{T} _ {g} (\rho (x)) u + \mathrm{T} _ {x} (\tau (g)) v).\tag{7}
\]

但是，根据假设 (b)，\(\mathrm{T}_{g}(\rho(x))\) 是单射，因而 \(\mathrm{T}_{(g,x)}(\theta)\) 是单射。其像是子空间 \(\mathrm{H}_{g,x}\) 与子空间的拓扑直和；前者由向量 \((v,\mathrm{T}_x(\tau(g))v)\)（其中 \(v\in \mathrm{T}_x(\mathrm{X})\)）构成，而后者为

\[
\mathrm{I} _ {g, x} = \{0 \} \times \mathrm{T} _ {g} (\rho (x)) (\mathrm{T} _ {g} (\mathrm{G})).
\]

根据假设 (b)，\(\mathrm{T}_g(\rho(x))(\mathrm{T}_g(\mathrm{G}))\) 有一个拓扑补 \(J_{g,x}\)（位于 \(\mathrm{T}_{gx}(\mathbf{X})\) 中）。因此 \(\mathrm{T}_{(g,x)}(\theta)\) 的像有拓扑补 \(\{0\} \times J_{g,x}\)。由此我们证明了 \(\theta\) 是从 \(\mathbf{G} \times \mathbf{X}\) 到 \(\mathbf{X} \times \mathbf{X}\) 的浸入。

由于 G 在 X 上自由作用，\(\theta\) 是单射。令 C 为 G 在 X 上定义的等价关系 R 的图。由于 G 作用适当，\(\theta\) 是从 \(G \times X\) 到 C 的同胚（General Topology, Chapter I, § 10, Proposition 2）。根据 Differentiable and Analytic Manifolds, R, 5.8.3，C 是 \(X \times X\) 的子流形，且 \(\theta\) 是从流形 \(G \times X\) 到流形 C 的同构。切空间 \(T_{(x,gx)}(C)\) 可与

\[
\mathrm{T} _ {(g, x)} (\theta) \left(\mathrm{T} _ {(g, x)} (\mathrm{G} \times \mathrm{X})\right) = \mathrm{H} _ {g, x} \oplus \mathrm{I} _ {g, x} \subset \mathrm{T} _ {(x, g x)} (\mathrm{X} \times \mathrm{X}).
\]

令 \(\mathrm{pr}_1\) 和 \(\mathrm{pr}_2\) 为 \(\mathbf{X} \times \mathbf{X}\) 到两个因子的典范投影。显然，\(\mathrm{T}_{(x,gx)}(\mathrm{pr}_1)\) 将 \(\mathrm{H}_{g,x}\) 映到 \(\mathrm{T}_x(\mathrm{X})\) 上，并且 \(\mathrm{T}_{(x,gx)}(\mathrm{pr}_1) \mid \mathrm{T}_{(x,gx)}(\mathrm{C})\) 的核是 \(\mathrm{I}_{g,x}\)。因此 \(\mathrm{pr}_1 \mid \mathrm{C}\) 是从 \(\mathrm{C}\) 到 \(\mathrm{X}\) 的浸没。根据 Differentiable and Analytic Manifolds, R, 5.9.5，R 是正规的。按定义，商集 X/G 上存在唯一的流形结构，使 \(\pi\) 是浸没。X/G 的底层拓扑是 X 的商拓扑（Differentiable and Analytic Manifolds, R, 5.9.4）。该拓扑是豪斯多夫的（General Topology, Chapter III, § 4, Proposition 3）。

对于所有 \(b \in \mathrm{X} / \mathrm{G}\)，存在一个开邻域 \(\mathrm{W}\)（它是 \(b\) 的邻域）以及映射 \(\sigma: \mathrm{W} \to \mathrm{X}\)，使 \(\pi \circ \sigma = \mathrm{id}_{\mathbf{W}}\)（Differentiable and Analytic Manifolds, R, 5.9.1）。令 \(\phi\) 为双射 \((g, w) \mapsto g\sigma(w)\)，它从 \(\mathrm{G} \times \mathrm{W}\) 到 \(\pi^{-1}(\mathrm{W})\)。它是 \(C^r\) 类的。于是 \(\pi(g\sigma(w)) = w\)，并且

\[
\theta^ {- 1} (\sigma (w), g \sigma (w)) = (g, \sigma (w))
\]

从而 \(\phi\) 的逆双射是 \(\mathbf{C}^r\) 类的。显然，\(\phi(gg', w) = g\phi(g', w)\) 对 \(w \in W\)、\(g \in G\)、\(g' \in G\) 成立。因此 \((\mathbf{X}, \mathbf{G}, \mathbf{X}/\mathbf{G}, \pi)\) 是主左纤维丛。

备注。在上述假设下，再设 H 是 \(C^{r}\) 类流形，且 \((x, h) \mapsto m(x, h)\) 是从 X × H 到 X 的 \(C^{r}\) 类映射，满足 \(m(gx, h) = gm(x, h)\) 对 \(x \in X, g \in G, h \in H\) 成立。令 n 为由 m 取商得到的从 \((\mathrm{X}/\mathrm{G}) \times \mathrm{H}\) 到 X/G 的映射。我们证明 n 是 \(C^{r}\) 类的。考虑图表

\[
\begin{array}{c} \mathrm{X} \times \mathrm{H} \xrightarrow {m} \mathrm{X} \\ \pi \times 1 \Biggl \downarrow \\ (\mathrm{X} / \mathrm{G}) \times \mathrm{H} \xrightarrow {n} \mathrm{X} / \mathrm{G} \end{array}
\]

该图是交换的，\(\pi \circ m\) 是 \(\mathbf{C}^r\) 类的，而 \(\pi \times 1\) 是满射浸没；于是只需应用 Differentiable and Analytic Manifolds, R, 5.9.5。

设 G 为李群，X 为 \(C^{r}\) 类流形，且 \((g, x) \mapsto xg\) 是 G 在 X 上的 \(C^{r}\) 类右作用律。令 \(\tau(g)x = \rho(x)g = xg\) 对 \(g \in G\)、\(x \in X\) 成立。这时

\[
\rho (x) = \rho (x g) \circ \gamma (g ^ {- 1}), \quad \rho (x) = \tau (g) \circ \rho (x) \circ \delta (g)\tag{3'}
\]

从而

(4')

\[
\mathrm{T} _ {g} (\rho (x)) = \mathrm{T} _ {e} (\rho (x g)) \circ \mathrm{T} _ {g} (\gamma (g ^ {- 1}))\tag{5'}
\]

\[
\mathrm{T} _ {g} (\rho (x)) = \mathrm{T} _ {x} (\tau (g)) \circ \mathrm{T} _ {e} (\rho (x)) \circ \mathrm{T} _ {g} (\delta (g)).
\]

另一方面，如果 \(\eta\) 表示映射 \((g, x) \mapsto xg\)，它从 \(G \times X\) 到 X，则公式 (6) 仍然成立。命题 9、其推论和命题 10 同样成立（在最后一个命题中，将 ``主左纤维丛'' 替换为 ``主右纤维丛''）。

